% Part: sets-functions-relations
% Chapter: relations
% Section: equivalence-relations

\documentclass[../../../include/open-logic-section]{subfiles}

\begin{document}

\olfileid{sfr}{rel}{eqv}

\olsection{Equivalentierelaties}

De identiteitsrelatie op een verzameling is reflexief, symmetrisch en
transitief. Veel relaties~$R$ hebben al deze drie eigenschappen.

\begin{defn}[Equivalentierelatie] 
Neem een relatie $R \subseteq A^2$. Als deze relatie reflexief, symmetrisch en
transitief is, noemen we haar een \emph{equivalentierelatie}.
!!^{element}s $x$ en $y$ uit~$A$ zijn \emph{$R$-equivalent} als~$Rxy$.
\end{defn}

Bij een equivalentierelatie hoort het begrip \emph{equivalentieklasse}.
Zo'n relatie verdeelt haar domein in losse delen. Binnen elk deel staan
alle objecten via de relatie met elkaar in verband. Objecten uit
verschillende delen staan niet zo met elkaar in verband. Soms willen
we gewoon over de delen \emph{zelf} praten. Daarvoor gebruiken we deze
definitie:

\begin{defn}\ollabel{def:equivalenceclass}
Laat $R \subseteq A^2$ een equivalentierelatie zijn. Neem een $x \in A$.
De \emph{equivalentieklasse} van $x$ in~$A$ is de verzameling
$\equivrep{x}{R} = \Setabs{y \in A}{Rxy}$. De verzameling van al deze
equivalentieklassen heet de \emph{quotiëntverzameling} die bij $A$ en~$R$ hoort:
$\equivclass{A}{R} = \Setabs{\equivrep{x}{R}}{x \in A}$. 
\end{defn}

De volgende stelling laat zien dat deze definitie werkt: de
equivalentieklassen vormen samen een partitie, dus een verdeling van~$A$:

\begin{prop}
Als $R \subseteq A^2$ een equivalentierelatie is, dan geldt $Rxy$ precies wanneer
$\equivrep{x}{R} = \equivrep{y}{R}$.
\end{prop}

\begin{proof}
Neem eerst aan dat $Rxy$ geldt. Kies een $z \in \equivrep{x}{R}$.
Volgens de definitie geldt dan $Rxz$. Omdat $R$ een equivalentierelatie
is, volgt $Ryz$. Dit gaat zo: uit $Rxy$ en de symmetrie van~$R$ volgt
$Ryx$. Uit dat omgekeerde verband, $Rxz$ en de transitiviteit van~$R$
volgt daarna~$Ryz$.
Dus $z \in \equivrep{y}{R}$. Daarom zit ieder element van de eerste
equivalentieklasse ook in de tweede, zodat $\equivrep{x}{R} \subseteq
\equivrep{y}{R}$. Hetzelfde argument, maar nu met de rollen omgedraaid,
geeft $\equivrep{y}{R} \subseteq \equivrep{x}{R}$. De twee verzamelingen
hebben dus precies dezelfde elementen. Volgens extensionaliteit is
$\equivrep{x}{R} = \equivrep{y}{R}$.

Neem nu aan dat $\equivrep{x}{R} = \equivrep{y}{R}$. Omdat $R$
reflexief is, geldt $Ryy$. Dus $y \in \equivrep{y}{R}$. Ook geldt
$y \in \equivrep{x}{R}$, want we nemen aan dat $\equivrep{x}{R} =
\equivrep{y}{R}$. Daarmee geldt $Rxy$.
\end{proof}

\begin{ex}
Een goed voorbeeld komt uit rekenen met resten, ook wel modulair rekenen
genoemd. Neem natuurlijke getallen $a$ en $b$ en een $n \in \PosInt$.
Dan betekent $a \equiv_n b$ dat
er bij deling van $a$ door~$n$ dezelfde rest overblijft als bij deling
van $b$ door~$n$. (Met symbolen staat er hetzelfde: $a \equiv_n b$ geldt
precies wanneer er een $k \in \Int$ is waarvoor $a - b = kn$.)
$\equiv_n$ is voor elk~$n$ een equivalentierelatie. Er zijn precies $n$
verschillende equivalentieklassen voor~$\equiv_n$. Anders gezegd:
$\equivclass{\Nat}{\equiv_n}$ heeft $n$ !!{element}s. Die klassen zijn:
de getallen die je zonder rest door $n$ kunt delen, oftewel
$\equivrep{0}{\equiv_n}$; de getallen waarbij na deling door $n$ rest~$1$
overblijft, oftewel $\equivrep{1}{\equiv_n}$; \ldots; en de getallen
waarbij na deling door~$n$ rest~$n-1$ overblijft, oftewel
~$\equivrep{n-1}{\equiv_n}$.
\end{ex}

\begin{prob}
Laat zien dat $\equiv_n$ voor elke $n \in \PosInt$ een
equivalentierelatie is en dat $\equivclass{\Nat}{\equiv_n}$ precies
$n$ elementen heeft.
\end{prob}
\end{document}
